% Chapter 4: Experimental Results
\chapter{Experimental Results}
\label{ch:results}

This chapter reports the experimental results of the MutBench framework presented in Chapter~\ref{ch:mutbench}.
Stage~1 (Sections~\ref{sec:variant_comparison}--\ref{sec:phylo_simulation}) conducts in-depth analysis on SARS-CoV-2, comparing five MutClust variants and two baseline methods through synthetic benchmark, sensitivity analysis, temporal generalization, multi-ground truth evaluation, and additional validation experiments.
Stage~2 (Section~\ref{sec:9pathogen_results} onwards) extends the evaluation to 9 RNA viruses with 9 scoring formulas $\times$ 39 detection algorithms, establishing the statistical evidence for pathogen-adaptive method selection.

%% ============================================================
\section{Single-Pathogen Benchmark: SARS-CoV-2}
\label{sec:single_pathogen_benchmark}

%% ------------------------------------------------------------
\subsection{Variant Comparison on Synthetic Data}
\label{sec:variant_comparison}

As a controlled signal recovery test, synthetic H-scores were generated for a SARS-CoV-2-like genome (29,903 positions) with 1,251 ground truth functional positions.
Table~\ref{tab:hotspot_scores} summarizes the results.

\begin{table}[htbp]
\centering
\caption{Synthetic benchmark results: method comparison with hotspot-score (100 subsamplings, 80\% retention)}
\label{tab:hotspot_scores}
\adjustbox{max width=\textwidth}{
\begin{tabular}{lrrrrr}
\toprule
\textbf{Method} & \textbf{Precision} & \textbf{Recall} & \textbf{F1} & \textbf{Stability} & \textbf{Hotspot-score} \\
\midrule
MutClust-Original & 0.991 & 0.504 & 0.668 & 0.418 & 0.638 \\
MutClust-Fixed & 0.988 & 0.506 & 0.669 & 0.445 & 0.646 \\
MutClust-dNdS & 0.988 & 0.506 & 0.669 & 0.445 & 0.646 \\
\textbf{MutClust-Hybrid} & \textbf{0.968} & \textbf{0.659} & \textbf{0.785} & \textbf{0.706} & \textbf{0.778} \\
HDBSCAN-Auto & 0.079 & 0.944 & 0.146 & 0.788 & 0.604 \\
\midrule
Frequency baseline & 0.404 & 0.965 & 0.569 & --- & --- \\
Sliding window & 0.157 & 0.974 & 0.270 & --- & --- \\
\bottomrule
\end{tabular}
}
\end{table}

MutClust-Hybrid achieves the best balance (F1 = 0.785, hotspot-score = 0.778) by combining CCM seed detection with HDBSCAN boundary determination. However, these synthetic results substantially overestimate real-world performance: on actual GISAID data (Section~\ref{sec:real_gisaid_benchmark}), F1 drops to 0.123 due to the extreme sparsity of real H-score distributions (342 of 29,903 positions nonzero). The synthetic benchmark validates the framework's ability to discriminate methods under controlled conditions, not absolute real-world performance.


%% ------------------------------------------------------------
\subsection{Multi-Ground Truth Evaluation}
\label{sec:multi_gt}

No single ground truth can capture all dimensions of mutational importance.
To address this, detection methods were evaluated against three complementary ground truths, each measuring a different biological phenomenon:

\begin{itemize}
\item \textbf{Functional regions} (439 AA positions, 34.5\% of Spike; this AA-level union of domain annotations includes broader boundary definitions than the nucleotide-level ground truth (1,251 nt, 417 AA after deduplication) used in the synthetic benchmark): Known domain annotations (RBD, NTD supersite, furin cleavage, fusion peptide, HR1, HR2).
\item \textbf{DMS-escape} (252 positions, 19.8\%): Experimentally measured single-mutation cell entry fitness effects~\cite{dadonaite2023dms} (top 20th percentile).
\item \textbf{Convergent evolution} (97 positions, 7.6\%): Literature-based sites under recurrent positive selection across independent SARS-CoV-2 lineages~\cite{cao2023ba286,carabelli2023convergent}. These include RBD immune evasion positions, ACE2 binding enhancement mutations, furin cleavage site adaptations, and S2 convergent substitutions (this expanded set includes all literature-reported convergent evolution positions, compared to the stricter 25-position Layer~A set used in the 9-pathogen benchmark).
\end{itemize}

The three ground truths are largely non-overlapping: the Jaccard similarity between DMS-escape and convergent evolution positions is 0.006, indicating near-complete independence.
This confirms that population frequency, experimental fitness, and evolutionary convergence capture fundamentally different aspects of mutational importance.

Critically, the best-performing method differed depending on the ground truth (Table~\ref{tab:multi_gt}).

\begin{table}[htbp]
\centering
\caption{Method comparison across multiple ground truths}
\label{tab:multi_gt}
\adjustbox{max width=\textwidth}{
\begin{tabular}{llrrr}
\toprule
\textbf{Ground Truth} & \textbf{Best method} & \textbf{F1} & \textbf{Enrichment} & \textbf{Positions} \\
\midrule
Functional regions (439) & TC-freq-v2 & 0.474 & --- & 439 \\
DMS-escape (252) & Saturated H-score & 0.315 & --- & 252 \\
Convergent evolution (97) & CH-score & 0.188 & 1.70$\times$ & 97 \\
Convergent evolution-strict$^*$ & TC-freq-v3 & 0.108 & 3.43$\times$ & --- \\
\bottomrule
\end{tabular}
}
\begin{minipage}{\textwidth}
\smallskip
\footnotesize $^*$Includes only strict convergent evolution sites confirmed from multiple sources. Jaccard(DMS, convergent evolution) = 0.006, confirming independence of the three ground truths.
\end{minipage}
\end{table}

For functional regions, TC-freq-v2 achieved the highest F1 (0.474); for DMS-escape, saturated H-score showed the best performance (F1 = 0.315); and for convergent evolution, CH-score achieved the highest F1 (0.188, enrichment ratio 1.70-fold).
For the strict convergent evolution subset (multi-source positions only), TC-freq-v3 achieved the highest enrichment ratio (3.43-fold, F1 = 0.108).

These differing rankings demonstrate that no single detection method is dominant across all biological dimensions of hotspot importance.
This multi-ground truth analysis provides a critical methodological insight: population frequency and experimental fitness are complementary rather than redundant measures of mutational importance.
Benchmarking frameworks that rely on a single ground truth risk systematic bias toward methods optimized for that specific biological dimension, potentially penalizing methods that capture orthogonal aspects of functional importance (Figure~\ref{fig:multi_ground_truth_top}, Figure~\ref{fig:multi_ground_truth_bottom}; Figure~\ref{fig:dms_evaluation_top}, Figure~\ref{fig:dms_evaluation_bottom}).

\begin{figure}[htbp]
\centering
\includegraphics[width=\textwidth]{multi_ground_truth_top.png}
\caption{Detection performance across multiple biological ground truths (Part 1). (A)~Best F1 by approach and ground truth, (B)~Enrichment ratio heatmap.}
\label{fig:multi_ground_truth_top}
\end{figure}

\begin{figure}[htbp]
\centering
\includegraphics[width=\textwidth]{multi_ground_truth_bottom.png}
\caption{Detection performance across multiple biological ground truths (Part 2). (C)~Ground truth overlap matrix (Jaccard similarity; diagonal shows set sizes), (D)~TC-freq-v3 score distributions for convergent vs non-convergent positions.}
\label{fig:multi_ground_truth_bottom}
\end{figure}

\begin{figure}[htbp]
\centering
\includegraphics[width=\textwidth]{dms_evaluation_top.png}
\caption{DMS-based ground truth evaluation (Part 1). (A)~F1 comparison between DMS-escape and functional-region ground truths, (B)~Cohen's d effect size for score separation across ground truths.}
\label{fig:dms_evaluation_top}
\end{figure}

\begin{figure}[htbp]
\centering
\includegraphics[width=\textwidth]{dms_evaluation_bottom.png}
\caption{DMS-based ground truth evaluation (Part 2). (C)~Scatter plot of approach scores vs DMS functional importance, (D)~Precision-recall curves for DMS-escape (top 20\%) ground truth.}
\label{fig:dms_evaluation_bottom}
\end{figure}

%% ------------------------------------------------------------
\subsection{Real GISAID H-score Benchmark}
\label{sec:real_gisaid_benchmark}

To address circularity concerns with synthetic data, a benchmark was conducted using the actual GISAID 2020--2021 H-score data used in the MutClust paper~\cite{youn2025mutclust} as input.
This data consists of H-scores calculated for each of 1,273 Spike amino acid positions from SARS-CoV-2 sequences collected from January 2020 to December 2021, representing actual observed data rather than synthetic generation.
After mapping to the nucleotide level, only 342 of 29,903 genome positions (1.14\%) had nonzero H-scores, exhibiting a sparse distribution qualitatively different from the uniform background noise of synthetic data.

The results are presented in Table~\ref{tab:real_gisaid}.

\begin{table}[htbp]
\centering
\caption{Method comparison on real GISAID 2020--2021 H-score data (342 nonzero positions, 1,251 ground truth nucleotide positions)}
\label{tab:real_gisaid}
\small
\adjustbox{max width=\textwidth}{
\begin{tabular}{lrrrrrr}
\toprule
\textbf{Method} & \textbf{Precision} & \textbf{Recall} & \textbf{F1} & \textbf{Enrichment} & \textbf{Stability} & \textbf{Hotspot-score} \\
\midrule
MutClust-Original & 0.333 & 0.012 & 0.023 & 7.97 & 0.372 & 0.239 \\
MutClust-Fixed & 0.295 & 0.079 & 0.125 & 7.04 & 0.784 & 0.386 \\
MutClust-dNdS & 0.295 & 0.079 & 0.125 & 7.04 & 0.784 & 0.386 \\
MutClust-Hybrid & 0.305 & 0.077 & 0.123 & 7.28 & 0.766 & 0.383 \\
HDBSCAN-Auto & --- & --- & --- & --- & --- & --- \\
\midrule
Frequency baseline & 0.289 & 0.079 & 0.124 & 6.92 & --- & --- \\
\textbf{Sliding window} & \textbf{0.323} & \textbf{0.600} & \textbf{0.420} & \textbf{7.73} & --- & --- \\
\bottomrule
\end{tabular}
}
\begin{minipage}{\textwidth}
\smallskip
\footnotesize HDBSCAN-Auto failed to form clusters due to insufficient nonzero positions (342). The sliding window achieved the highest F1 (0.420) through its wide detection window (20 positions).
\end{minipage}
\end{table}

Three key differences from the synthetic benchmark were observed.
First, F1 of all MutClust variants dropped substantially (e.g., MutClust-Hybrid: 0.785 $\to$ 0.123) due to sparsity: many ground truth positions (conserved regions HR1, HR2) have H-score = 0 and cannot be detected by density-based methods.
Second, method rankings shifted---MutClust-Hybrid's advantage over MutClust-Fixed disappeared (0.383 vs.\ 0.386), and the sliding window achieved the highest F1 (0.420) because wide detection windows capture adjacent functional regions in sparse data.
Third, enrichment ratios remained consistently high at 6.9--8.0-fold, confirming biological validity despite low absolute F1.
These results reinforce a core lesson: method superiority depends on H-score distribution characteristics, and synthetic conclusions should not be unconditionally extrapolated to real data.

%% ------------------------------------------------------------
\subsection{Comparison with External Methods}
\label{sec:external_comparison}

Two independent methods---OPTICS (density-based clustering, min\_samples=5, $\xi$=0.05) and KDE (kernel density peak detection, 90th percentile threshold)---were evaluated on the same synthetic benchmark (Table~\ref{tab:external_methods}).

\begin{table}[htbp]
\centering
\caption{Benchmark results including external independent methods. Synthetic H-score data, 7 Spike functional region ground truth.}
\label{tab:external_methods}
\small
\adjustbox{max width=\textwidth}{
\begin{tabular}{lcccccc}
\toprule
Method & Precision & Recall & F1 & Stability & Hotspot-score \\
\midrule
MutClust-Hybrid & 0.968 & 0.659 & 0.785 & 0.706 & \textbf{0.778} \\
MutClust-Original & 0.991 & 0.504 & 0.668 & 0.418 & 0.638 \\
MutClust-Fixed & 0.988 & 0.506 & 0.669 & 0.445 & 0.646 \\
MutClust-dNdS & 0.988 & 0.506 & 0.669 & 0.445 & 0.646 \\
HDBSCAN-Auto & 0.079 & 0.944 & 0.146 & 0.788 & 0.604 \\
KDE & \textbf{1.000} & 0.049 & 0.093 & 0.622 & 0.557 \\
OPTICS & 0.087 & 0.546 & 0.150 & 0.473 & 0.369 \\
\bottomrule
\end{tabular}
}
\end{table}

MutClust-Hybrid maintained the highest hotspot-score (0.778) even against external methods. Parameter sensitivity analysis across 30 OPTICS and 36 KDE configurations confirmed this superiority is structural rather than parameter-dependent: even at OPTICS's best parameters (F1=0.707), precision (0.563) remained far below MutClust-Hybrid (0.968).

\textbf{Stage 1 summary.} MutClust-Hybrid achieves the best hotspot-score (0.778) on synthetic data, but all methods show reduced performance on real GISAID data (F1: 0.785 $\to$ 0.123) due to sparse H-score distributions, while maintaining biological validity (enrichment $>$ 6.9-fold).

%% ============================================================
\section{Robustness and Sensitivity Validation}
\label{sec:robustness_validation}

%% ------------------------------------------------------------
\subsection{Statistical Validation}
\label{sec:statistical_validation}

To evaluate the statistical robustness of the benchmark results, bootstrap confidence interval estimation, permutation tests (random label shuffling + circular permutation), pairwise comparison tests, and Fisher's exact test were performed.
All bootstrap confidence intervals reported in this section were estimated using the bias-corrected and accelerated (BCa) method~\cite{efron1987better} to correct for bias and asymmetry in small samples ($n = 7$).

\textbf{Permutation tests and hotspot-score null distribution.}
Both random label shuffling (1,000 iterations, all methods $p < 0.001$) and circular permutation tests (200 iterations, preserving spatial autocorrelation) confirmed that all detection results are statistically significant.
Under circular permutation, all 5 methods showed F1 values substantially exceeding the null distribution ($z > 9.2$, $p < 0.005$ for all; e.g., MutClust-Hybrid: F1 = 0.785, null mean = 0.023).
The composite hotspot-scores were similarly significant: MutClust-Hybrid (HS = 0.778, null mean = 0.251, $p < 0.005$), MutClust-Fixed (HS = 0.646, null = 0.162), HDBSCAN-Auto (HS = 0.604, null = 0.441).
These results confirm that detected hotspot--functional region overlap reflects genuine biological signal rather than chance or spatial autocorrelation artifacts.

\textbf{Fisher's exact test.}
The statistical significance of enrichment ratios reported in the multi-ground truth evaluation (Section~\ref{sec:multi_gt}) was validated using Fisher's exact test.
After applying Benjamini-Hochberg FDR correction ($q < 0.05$) to 128 comparisons, 3 comparisons were significant:
functional region enrichment of TC-freq-v2 (1.16-fold, $p_{\text{adj}}$ = 0.002),
functional region enrichment of TC-freq-v3 (1.59-fold, $p_{\text{adj}}$ = 0.023),
convergent evolution enrichment of CH-score (1.70-fold, $p_{\text{adj}}$ = 0.023).
Additionally, 2 comparisons were significant at $q < 0.10$ (CH-score convergent evolution top10pct, E-score convergent evolution top20pct).
The enrichment of TC-freq-v3 in the strict convergent evolution subset (3.43-fold, $p$ = 0.013, $p_{\text{adj}}$ = 0.161) and the convergent evolution enrichment of saturated H-score (1.85-fold, $p$ = 0.048, $p_{\text{adj}}$ = 0.269) did not reach significance after BH FDR correction; however, the enrichment trend at convergent evolution positions was consistently observed across multiple approaches (CH-score, E-score, TC-freq-v3).

\textbf{BCa bootstrap confidence intervals.}
Region-level BCa bootstrap (1,000 iterations, $n = 7$ functional regions)~\cite{efron1987better} confidence intervals for hotspot-score are presented in Table~\ref{tab:bootstrap_ci}.

\begin{table}[htbp]
\centering
\caption{Region-level BCa bootstrap 95\% confidence intervals for hotspot-score (1,000 resamplings of 7 functional regions, bias-corrected and accelerated)}
\label{tab:bootstrap_ci}
\adjustbox{max width=\textwidth}{
\begin{tabular}{lrrrcc}
\toprule
\textbf{Method} & \textbf{HS$_{\text{region}}$} & \textbf{BCa lower} & \textbf{BCa upper} & \textbf{$z_0$} & \textbf{$a$} \\
\midrule
\textbf{MutClust-Hybrid} & \textbf{0.476} & \textbf{0.323} & \textbf{0.635} & $-$0.090 & 0.010 \\
HDBSCAN-Auto & 0.468 & 0.339 & 0.567 & $-$0.003 & $-$0.025 \\
MutClust-Fixed & 0.311 & 0.196 & 0.534 & 0.030 & 0.070 \\
MutClust-Original & 0.302 & 0.187 & 0.505 & $-$0.015 & 0.070 \\
\bottomrule
\end{tabular}
}
\begin{minipage}{\textwidth}
\smallskip
\footnotesize HS$_{\text{region}}$: Mean of hotspot-scores across 7 regions. Unlike the aggregate hotspot-score (Table~\ref{tab:hotspot_scores}), the region-level mean assigns equal weight to high-signal and low-signal regions, resulting in lower absolute values.
\end{minipage}
\end{table}

The region-level BCa bootstrap quantifies uncertainty in ground truth definition by resampling the 7 functional regions.
CCM-based methods show right-skewed distributions ($a = 0.070$) driven by RBD inclusion, while HDBSCAN-Auto's broad detection yields a narrower, symmetric interval (0.339--0.567); all point estimates fall within the 95\% BCa CI, confirming statistical validity.

\textbf{Hotspot-score weight sensitivity.}
\phantomsection\label{tab:weight_sensitivity}
MutClust-Hybrid maintained rank 1 in 71.9\% (123/171) of weight combinations, switching to HDBSCAN-Auto only under extreme recall emphasis (weight $\geq$ 0.60).
Pairwise bootstrap comparisons confirmed all score differences are significant ($p < 0.001$).

%% ------------------------------------------------------------
\subsection{Sensitivity Analysis}
\label{sec:sensitivity_analysis}

\subsubsection{Sample Size Sensitivity}

To verify that the benchmark results are not artifacts of a specific sample size, the effect of varying the number of input sequences on detection performance was analyzed.
To determine the minimum number of input sequences required for stable hotspot detection, MutClust-Hybrid was evaluated across sample sizes ranging from 100 to 10,000 sequences (5 replicates per size).

H-score rank correlation (Spearman $\rho$) between subsamples and the full dataset exceeded 0.997 at $n = 1{,}000$ and reached 1.000 at $n = 10{,}000$.
Hotspot-score stabilized at approximately 0.62 by $n = 1{,}000$, with minimal inter-replicate variance (standard deviation $< 0.015$).
Even at the smallest sample size ($n = 100$), hotspot-score remained above 0.52 with Spearman $\rho > 0.94$, indicating that the rank order of mutation intensity is preserved even in limited datasets (Figure~\ref{fig:sample_size_convergence}).

These results suggest that approximately 1,000 sequences are sufficient for stable detection results, and the 10,000 sequences in the benchmark dataset provide considerable margin above this threshold.

\begin{figure}[htbp]
\centering
\includegraphics[width=0.8\textwidth]{sample_size_convergence.png}
\caption{Sample size convergence of hotspot metrics. Stable detection performance (Hotspot-score $\approx$ 0.62) is achieved starting at approximately 1,000 sequences.}
\label{fig:sample_size_convergence}
\end{figure}

\subsubsection{Parameter Sensitivity}

The sensitivity of detection performance to MutClust parameters is visualized in Figure~\ref{fig:mutbench_sensitivity}.

\begin{figure}[htbp]
\centering
\includegraphics[width=0.9\textwidth]{sensitivity_heatmap.png}
\caption{MutClust parameter sensitivity heatmap. Full-factorial grid search over 64 parameter combinations ($\gamma \times d \times$ MinPts) shows that F1 performance is most sensitive to $\gamma$ (epsilon scaling factor), followed by $d$ (decay factor), with minimal influence of MinPts.}
\label{fig:mutbench_sensitivity}
\end{figure}

\phantomsection\label{tab:hdbscan_sensitivity}
A grid search over 4 HDBSCAN parameter configurations (min\_cluster\_size $\in \{5,10,15\}$, min\_samples $\in \{2,3,5,7\}$) showed minimal impact on detection results (F1 range: 0.143--0.150), confirming robustness to parameter choices.
This parameter insensitivity is a key advantage of density-based approaches and provides justification for MutClust-Hybrid's use of HDBSCAN for boundary determination.

\subsubsection{dN/dS Threshold Sensitivity}
\phantomsection\label{tab:dnds_threshold}

dN/dS threshold variation (0.5--5.0) did not alter the ranking of top methods (MutClust-Hybrid highest across all thresholds, F1 stable at 0.774--0.788 for thresholds 0.5--1.0), confirming robustness to this parameter choice.

%% ------------------------------------------------------------
\subsection{Temporal and Structural Robustness}
\label{sec:temporal}

\textbf{Temporal generalization (GISAID 2024--2025).}
To evaluate temporal generalizability, 9,957 SARS-CoV-2 sequences from January 2024 to December 2025 were analyzed.
The validation revealed \textbf{H-score saturation}: the 2024--2025 dataset produced nonzero H-scores at all 1,273 Spike positions (mean = 8.47, SD = 0.098), compared to only 114 positions in the 2020--2021 dataset.
This nearly uniform distribution eliminates discriminative power, causing all methods to converge to the ground truth base rate precision ($\sim$0.33) (Table~\ref{tab:temporal}).
This saturation reflects a biological reality---Omicron-era lineages have accumulated extensive mutations relative to Wuhan-Hu-1---and validates the critical limitation that H-scores computed against a fixed ancestral reference lose discriminative power as the viral population diverges.

\begin{table}[htbp]
\centering
\caption{Temporal validation results using GISAID 2024--2025 data}
\label{tab:temporal}
\adjustbox{max width=\textwidth}{
\begin{tabular}{lrrrr}
\toprule
\textbf{Method} & \textbf{Precision} & \textbf{Recall} & \textbf{Detections} & \textbf{Clusters} \\
\midrule
MutClust-Fixed & 0.328 & 1.000 & 3,819 & 1 \\
MutClust-dNdS & 0.328 & 1.000 & 3,819 & 1 \\
MutClust-Hybrid & $\sim$0.33 & $\sim$1.0 & 3,225 & 377 \\
\bottomrule
\end{tabular}
}
\begin{minipage}{\textwidth}
\smallskip
\footnotesize Nonzero H-scores were observed at all 1,273 Spike positions, eliminating discriminative power. Precision converges to the ground truth base rate ($\sim$0.33).
\end{minipage}
\end{table}

\textbf{Indel exclusion impact.}
\label{sec:indel_impact}
H-score is a fixed-coordinate metric that cannot capture insertions or deletions.
Analysis of the same 2024--2025 sequences confirmed that 99.8\% carry at least one functional indel (e.g., del69-70, Y144del), yet all indel positions have H-score = 0.000.
Although the quantitative F1 impact is negligible ($\Delta$F1 = $-$0.002, as 5 indel positions add only 27 nt to the 1,251 nt ground truth), indel exclusion represents a structural limitation: biologically important structural variants are systematically invisible to position-level scoring.
The implications are discussed in Chapter~\ref{ch:conclusion}.

%% ============================================================
\section{Cross-Pathogen and Structural Validation}
\label{sec:cross_pathogen_structural}

%% ------------------------------------------------------------
\subsection{Influenza A H3N2 Cross-Validation}
\label{sec:influenza_validation}

To evaluate the cross-pathogen generalizability of MutBench, the same benchmark pipeline was applied to the hemagglutinin (HA) gene of Influenza A H3N2.
A total of 1,967 human-derived H3N2 HA protein sequences (540--580 AA) were collected from the NCBI Influenza Virus Resource, and H-scores were calculated from position-wise mutation frequency and Shannon entropy relative to the consensus sequence.

As ground truth, the 5 antigenic sites (antigenic sites A--E) of H3N2 HA were used (Table~\ref{tab:influenza})~\cite{wiley1981structural,wilson1990structural}, which are established functional regions defined by decades of serological and structural studies (101 AA positions, 303 nucleotide positions). Smith et al.~\cite{smith2004antigenic} demonstrated through antigenic cartography that a small number of substitutions at these sites drive major antigenic transitions in H3N2 evolution.

\begin{table}[htbp]
\centering
\caption{Influenza A H3N2 HA cross-pathogen benchmark results (1,967 sequences, 5 antigenic site ground truth)}
\label{tab:influenza}
\small
\adjustbox{max width=\textwidth}{
\begin{tabular}{lrrrrrr}
\toprule
\textbf{Method} & \textbf{Precision} & \textbf{Recall} & \textbf{F1} & \textbf{Enrichment} & \textbf{Stability} & \textbf{Hotspot-score} \\
\midrule
MutClust-Original & 0.182 & 1.000 & 0.309 & 1.02 & 0.535 & 0.572 \\
\textbf{MutClust-Fixed} & \textbf{0.183} & \textbf{1.000} & \textbf{0.309} & \textbf{1.02} & \textbf{0.799} & \textbf{0.661} \\
MutClust-Hybrid & 0.178 & 0.812 & 0.292 & 1.00 & 0.742 & 0.577 \\
HDBSCAN-Auto & 0.155 & 0.366 & 0.218 & 0.87 & 0.750 & 0.424 \\
KDE & 0.410 & 0.083 & 0.138 & 2.30 & 0.605 & 0.366 \\
\midrule
Frequency baseline & 0.421 & 0.238 & 0.304 & 2.36 & --- & --- \\
\textbf{Sliding window} & \textbf{0.402} & \textbf{0.462} & \textbf{0.430} & \textbf{2.25} & --- & --- \\
\bottomrule
\end{tabular}
}
\end{table}

Two key observations emerge. First, method rankings reversed: MutClust-Hybrid (hotspot-score 0.778 on SARS-CoV-2) was outperformed by MutClust-Fixed (0.661 vs.\ 0.577) on H3N2, because all 566 HA positions have nonzero H-scores, eliminating HDBSCAN's advantage.
Second, enrichment ratios decreased substantially (0.87--2.36-fold vs.\ 1.9--23.7-fold for SARS-CoV-2), reflecting the short HA gene length and high background mutation rate (Table~\ref{tab:cross_pathogen}, Figure~\ref{fig:cross_pathogen_comparison}).

\begin{table}[htbp]
\centering
\caption{Two-pathogen cross-comparison: pathogen dependence of the best-performing method}
\label{tab:cross_pathogen}
\small
\adjustbox{max width=\textwidth}{
\begin{tabular}{llrr}
\toprule
& & \multicolumn{2}{c}{\textbf{Hotspot-score}} \\
\cmidrule(lr){3-4}
\textbf{Method} & \textbf{Type} & \textbf{SARS-CoV-2} & \textbf{H3N2} \\
\midrule
MutClust-Original & CCM & 0.638 & 0.572 \\
MutClust-Fixed & CCM & 0.646 & \textbf{0.661} \\
MutClust-Hybrid & CCM+HDBSCAN & \textbf{0.778} & 0.577 \\
HDBSCAN-Auto & Density-based & 0.604 & 0.424 \\
\bottomrule
\end{tabular}
}
\end{table}

\begin{figure}[htbp]
\centering
\includegraphics[width=\textwidth]{cross_pathogen_comparison.png}
\caption{Cross-pathogen comparison of detection method rankings between SARS-CoV-2 and H3N2. Method rankings differ substantially across pathogens, motivating pathogen-adaptive selection.}
\label{fig:cross_pathogen_comparison}
\end{figure}

%% ------------------------------------------------------------
\subsection{3D Structure Analysis}
\label{sec:structure_analysis}

To validate the structural plausibility of linear sequence-based hotspot detection, the 3D spatial distribution of detected hotspot positions was analyzed using C$\alpha$ coordinates from PDB 6VSB (Spike protein prefusion structure, Wrapp et al. 2020).
The proportion of hotspot position pairs within 8\,\AA\ contact distance was 5.80\%, showing a \textbf{5.92-fold higher spatial contact enrichment} compared to random position pairs (0.98\%) ($z$ = 13.34, based on 1,000 random permutation tests).
This quantitatively demonstrates that hotspots detected from linear sequences are also spatially clustered in 3D structure, supporting the biological validity of the detection results.

%% ------------------------------------------------------------
\subsection{Phylogenetic Non-Independence Simulation}
\label{sec:phylo_simulation}

To quantitatively evaluate whether H-scores can be biased by phylogenetic structure, coalescent-based simulations and real-data validation were performed.

\subsubsection{Simulation and ML Validation}
\label{subsec:phylo_correction_poc}

In a coalescent simulation (100 replicates, $N = 500$), H-score ranked a single-origin founder mutation (frequency 90\%, D614G-like) above convergent mutations (15 independent origins, E484K-like) in 100\% of replicates, whereas Fitch parsimony-based homoplasy counts correctly ranked convergent mutations at the top in 100\% ($\rho = -0.876$, $p = 1.28 \times 10^{-64}$ between H-score and independent origin count).
\phantomsection\label{subsec:treetime_ml_validation}An extended validation using TreeTime~\cite{sagulenko2018treetime} ML-based ancestral reconstruction ($N = 2{,}000$, 50 replicates) confirmed that both Fitch parsimony and TreeTime-ML achieved AUC = 1.000 for discriminating convergent from non-convergent mutations, while H-score achieved AUC = 0.000.
In the MutBench ground truth, D614G is deliberately excluded from Layer~A and Layer~B because it reached fixation via a single founder event~\cite{korber2020tracking}.
These results demonstrate that \textbf{frequency-based scores systematically conflate founder effects with positive selection}, though this simulation assumes ideal conditions; at GISAID scale ($> 16$M sequences), polytomies, computational cost, and recombination introduce additional challenges.

\subsubsection{Real GISAID Validation}
\label{subsec:treetime_gisaid_validation}

Validation on 500 GISAID sequences (2024--2025) with TreeTime ML ancestral reconstruction confirmed a \textbf{weak negative correlation} between H-score and homoplasy score ($\rho = -0.107$, $p = 1.36 \times 10^{-4}$ across 1,273 Spike positions), demonstrating that the two scoring systems capture fundamentally different evolutionary signals.
Key positions illustrate the divergence: P681 ranks 2nd by H-score but 1,134th by homoplasy (founder effect), while L452 and K417 have H-score = 0 but homoplasy of 91 (convergent evolution).

\subsubsection{Homoplasy-Based Scoring Sensitivity Analysis}
\label{subsec:homoplasy_sensitivity}

To validate with real phylogenetic data, TreeTime ML ancestral reconstruction was applied to 500 SARS-CoV-2 Spike sequences (2020-H1, NCBI); per-position homoplasy counts were computed and evaluated through seven detection methods against the Layer~A ground truth (25 convergent evolution positions).

\begin{table}[htbp]
\centering
\caption{Comparison of detection method MCC under H-score versus TreeTime ML-inferred homoplasy counts for SARS-CoV-2 Spike (500 sequences, 2020-H1).}
\label{tab:phylo_sensitivity}
\small
\begin{tabular}{lccccc}
\toprule
Detector & \multicolumn{2}{c}{H-score} & \multicolumn{2}{c}{TreeTime Homoplasy} & $\Delta$MCC \\
\cmidrule(lr){2-3} \cmidrule(lr){4-5}
 & Best formula & MCC & MCC & Precision & \\
\midrule
FreqThresh(p=85)     & rank(P*E) & 0.015 & $-0.036$ & 0.000 & $-0.051$ \\
FreqThresh(p=90)     & P*E       & 0.047 & $-0.029$ & 0.000 & $-0.076$ \\
FreqThresh(p=95)     & E*rare    & 0.076 & $-0.015$ & 0.000 & $-0.091$ \\
AdaptiveKnee         & E*rare    & 0.116 & $-0.033$ & 0.000 & $-0.149$ \\
LocalContrast(z=1.5) & minority\_E & 0.068 & $-0.033$ & 0.000 & $-0.101$ \\
Wavelet(t=1.5)       & E*rare    & 0.086 & $-0.016$ & 0.000 & $-0.102$ \\
ScoreDBSCAN(ep=10)   & rank(P*E) & 0.023 & $+0.033$ & 0.025 & $+0.010$ \\
\midrule
Mean                 &           & 0.062 & $-0.018$ &       & $-0.080$ \\
\bottomrule
\end{tabular}
\end{table}

All detectors except ScoreDBSCAN yield negative MCC under homoplasy scoring (mean $-0.018$), primarily due to temporal mismatch: the 2020-H1 sequences predate Variants of Concern emergence, so Layer~A convergent evolution positions had not yet arisen as recurrent mutations.
The founder-effect position D614G correctly received a homoplasy count of 0, confirming that TreeTime distinguishes founder-driven fixation from convergent evolution.
The Spearman rank correlation between method rankings under H-score and TreeTime homoplasy is $\rho = 0.036$ ($p = 0.939$), demonstrating that method rankings are conditional on the scoring system and temporal scope; the implications for adaptive method selection are discussed in Chapter~\ref{ch:discussion}.

\textbf{Cross-pathogen summary.} Method rankings differ between SARS-CoV-2 and H3N2 (MutClust-Fixed outperforms MutClust-Hybrid on H3N2), and phylogenetic non-independence analysis reveals that H-score conflates founder effects with positive selection ($\rho = -0.876$). These findings motivate the large-scale 9-pathogen benchmark.

%% ============================================================
\section{9-Pathogen Large-Scale Benchmark}
\label{sec:9pathogen_results}

\vspace{1em}
\noindent\rule{\textwidth}{0.4pt}
\vspace{0.5em}

Building on the Stage~1 findings, a Stage~2 benchmark was conducted to systematically evaluate 39 detection methods from 14 algorithm families combined with 9 scoring formulas across 9 pathogens (Section~\ref{sec:9pathogen_benchmark}).
Stage~2 uses MCC rather than hotspot-score because the stability component (Jaccard-based subsampling reproducibility) requires bootstrap validation that has not been cross-validated across all 9 pathogens with their widely varying alignment lengths (340--1,330 AA) and MSA sizes (530--5,019 unique sequences). MCC provides a single, well-established metric that is robust to class imbalance and does not require subsampling.

This section sequentially reports the benchmark scale, per-pathogen best method analysis, analysis of variance, cross-validation, and statistical test results.

%% ------------------------------------------------------------
\subsection{Benchmark Scale and Data}
\label{subsec:9pathogen_scale}

This section presents the full-scale 9-pathogen benchmark results that constitute the primary evidence for the dissertation's second contribution (statistical demonstration of pathogen-adaptive method selection necessity).

A total of 3,159 evaluations (9 scoring formulas $\times$ 39 detection algorithms $\times$ 9 pathogens) were conducted.
Of these, 741 (23.5\%) evaluations resulted in the detection method identifying 0 positions (n\_detected = 0), because certain SWAN and ScoreDBSCAN variants failed to find positions satisfying the $z$-normalization threshold or density criteria in the score distributions of specific pathogens.
These zero-detection evaluations are treated as MCC = 0. However, the ANOVA analysis used 2,418 valid evaluations excluding zero-detection cases, with Type III sum of squares applied.
Each evaluation computed MCC, F1, precision, recall, enrichment ratio, and constrained region FPR against the 3-layer ground truth (Layer A: convergent evolution positions, Layer B: constrained regions, Layer C: DMS).
DMS F1 was computed only for the 4 pathogens (SARS-CoV-2, H3N2, HIV-1, RSV) where experimental fitness data are available.
All FASTA sequences were collected from NCBI GenBank using actual accession numbers.

%% ------------------------------------------------------------
\subsection{Per-Pathogen Best Method Analysis}
\label{subsec:9pathogen_best_combos}

The central question of Stage 2 is whether the same scoring--detection combination performs best across all pathogens, or whether each pathogen requires a different optimal combination.
The best scoring--detection combination was selected for each of the 9 pathogens based on MCC, and \textbf{all 9/9 pathogens exhibited distinct best combinations} (Table~\ref{tab:9pathogen_best}).

\begin{table}[htbp]
\centering
\caption{Best scoring--detection combination for each of the 9 pathogens (by MCC)}
\label{tab:9pathogen_best}
\small
\adjustbox{max width=\textwidth}{
\begin{tabular}{llll}
\toprule
\textbf{Pathogen} & \textbf{Best scoring} & \textbf{Best detection} & \textbf{MCC} \\
\midrule
SARS-CoV-2 & E*rare & CUSUM ($\delta$=0.5, $h$=4.0) & 0.118 \\
H3N2 & E\_only & SlidingTest ($p$=0.05) & 0.549 \\
Norovirus & P*E*rare & GradPeak ($p$=0.5) & 0.438 \\
HIV-1 & E*rare & SlidingTest ($p$=0.01) & 0.182 \\
Dengue & E*rare & FreqThresh ($p$=80) & 0.256 \\
RSV & P*E & FreqThresh ($p$=95) & 0.249 \\
Influenza B & H-score & Wavelet ($t$=2.0) & 0.361 \\
MERS & rank(P*E) & SlidingTest ($p$=0.05) & 0.088 \\
HCV & E\_only & CUSUM ($\delta$=0.5, $h$=5.0) & 0.443 \\
\bottomrule
\end{tabular}
}
\begin{minipage}{\textwidth}
\smallskip
\footnotesize Nine unique (scoring, detection) combinations were selected as best-performing across the 9 pathogens, with no single combination being best for two or more pathogens simultaneously.
\end{minipage}
\end{table}

\begin{figure}[htbp]
\centering
\includegraphics[width=\textwidth]{scoring_detection_heatmap.png}
\caption{Scoring--detection performance heatmap across 9 pathogens. Each cell represents mean MCC for a scoring--detection family combination. The heterogeneous pattern confirms that no single combination dominates across all pathogens. Detailed MCC values for all scoring--detection combinations are available in Table~\ref{tab:9pathogen_best}; the heatmap provides a visual overview of performance patterns.}
\label{fig:scoring_detection_heatmap}
\end{figure}

The key observations from these results are as follows.
First, the best scoring formula differs across pathogens: E*rare is the best-performing formula for SARS-CoV-2, HIV-1, and Dengue, but E\_only is best for H3N2 and HCV, H-score for Influenza B, P*E for RSV, and rank(P*E) for MERS.
Second, the best detection family also differs across pathogens: SlidingTest (3 occurrences), CUSUM (2), FreqThresh (2), Wavelet (1), and GradPeak (1) are distributed.
Third, the MCC range spans from 0.088 (MERS) to 0.549 (H3N2), exhibiting large variation that reflects structural differences in mutation patterns across pathogens (entropy distribution, proportion of conserved regions, number of positively selected sites).

%% ------------------------------------------------------------
\subsection{Statistical Analysis}
\label{subsec:9pathogen_statistical}

To determine whether the observed method--pathogen dependence is statistically significant, three complementary statistical analyses were performed.

\textbf{Three-way ANOVA ($\omega^2$).}
\label{subsec:9pathogen_anova}
To identify the major sources of performance variation, a three-way ANOVA was performed with MCC as the dependent variable.
The results using the bias-corrected effect size $\omega^2$ (Eq.~\ref{eq:omega_squared}) are presented in Table~\ref{tab:9pathogen_anova}.

\begin{table}[htbp]
\centering
\caption{Three-way ANOVA results: $\omega^2$ effect size of each factor on MCC}
\label{tab:9pathogen_anova}
\adjustbox{max width=\textwidth}{
\begin{tabular}{llr}
\toprule
\textbf{Factor} & \textbf{Type} & \textbf{$\omega^2$} \\
\midrule
pathogen & Main effect & 0.260 \\
family $\times$ pathogen & Two-way interaction & \textbf{0.285} \\
score $\times$ pathogen & Two-way interaction & 0.083 \\
score & Main effect & 0.027 \\
family & Main effect & 0.025 \\
score $\times$ family & Two-way interaction & 0.013 \\
\midrule
Residual & --- & 0.307 \\
\bottomrule
\end{tabular}
}
\begin{minipage}{\textwidth}
\smallskip
\footnotesize family: detection method family (14 levels), score: scoring formula (9 levels), pathogen: pathogen (9 levels). $\omega^2$ is the bias-corrected estimate of the proportion of total variance explained by each factor.
\end{minipage}
\end{table}

\begin{figure}[htbp]
\centering
\includegraphics[width=\textwidth]{variance_decomposition.png}
\caption{Variance decomposition of hotspot detection performance. The detection family $\times$ pathogen interaction ($\omega^2 = 0.285$) is the largest source of variance, exceeding all main effects.}
\label{fig:variance_decomposition}
\end{figure}

\textbf{Key finding}: The detection family $\times$ pathogen interaction ($\omega^2 = 0.285$) is the \textbf{largest source of variance, exceeding all main effects and other interactions}. Under Cohen's effect size benchmarks~\cite{cohen1988statistical}, $\omega^2 > 0.14$ constitutes a large effect, confirming that this interaction is substantively large.
This means that the best-performing detection method varies systematically depending on the pathogen, and no single pathogen-independent best method exists.

This finding directly resolves the third problem from Chapter~\ref{ch:introduction}: the interaction dominance ($\omega^2 = 0.285$) provides definitive statistical proof that detection performance is fundamentally pathogen-dependent. The pathogen main effect ($\omega^2 = 0.260$) is also large, while the main effects of scoring ($\omega^2 = 0.027$) and detection family ($0.025$) are small, confirming that ``which method is applied to which pathogen'' is the dominant factor.
The residual ($1 - 0.693 = 30.7\%$) reflects within-cell variability.
The three-way interaction (score $\times$ family $\times$ pathogen) was not estimated due to insufficient degrees of freedom from the unbalanced design (33--38 valid methods per pathogen); since the two-way interaction already dominates, this omission does not affect the core conclusions.

\textbf{LOPO Cross-Validation.}
\label{subsec:9pathogen_lopo}
LOPO (Leave-One-Pathogen-Out) cross-validation was performed to evaluate whether the best combination selected from 8 training pathogens is also best-performing for the excluded test pathogen.

\begin{itemize}
\item \textbf{Combination match rate}: 0/9 (0\%). In all 9 iterations, the training-best combination and the test oracle combination were mismatched.
\item \textbf{Mean generalization ratio}: 4.9\%. The test MCC of the training-best combination was only 4.9\% of the oracle MCC on average.
\item \textbf{Mean performance gap}: 0.277. The average MCC difference between the training-best combination and the oracle was 0.277, representing a substantial loss of oracle performance.
\item \textbf{Most frequent training-best combination}: E\_only + SWAN was selected in 8/9 iterations, an artifact of SWAN producing valid results only for Norovirus (MCC = 0.311) and MCC = 0.000 elsewhere.
\end{itemize}

SWAN achieves non-zero detection for only 1 of 9 pathogens (Norovirus, MCC = 0.311), yet appears as the training-best detector in 8/9 LOPO iterations. Excluding SWAN improves the mean generalization ratio from 4.9\% to 11.6\%, but the combination match rate remains 0/9, confirming that non-transferability does not depend on the SWAN artifact.

These results demonstrate that the 0/9 match rate is a \textbf{key finding}: the premise of a single best-performing method does not hold, providing the strongest statistical evidence for the necessity of pathogen-adaptive method selection.

\textbf{Friedman Test.}
\label{subsec:9pathogen_friedman}
Using the 9 pathogens as blocks, a Friedman test was performed on the performance differences among the top 20 scoring--detection combinations.
$k = 20$ corresponds to approximately the top 7\% of the total valid combinations (247--305 per pathogen) and represents the competitive combination group while maintaining a treatment count that is not excessive relative to the number of blocks ($b = 9$)\footnote{To verify the robustness of the $k$ selection, the Friedman test was also performed at $k = 10$, $k = 30$, and $k = 50$. $k = 10$: $\chi^2 = 22.1$, $p = 0.009$; $k = 30$: $\chi^2 = 58.7$, $p = 0.001$; $k = 50$: $\chi^2 = 87.3$, $p < 0.001$. All $k$ values yielded $p < 0.01$, confirming that the significance of performance differences among top combinations is robust to the choice of $k$.}.

\begin{equation}
\chi^2 = 42.44, \quad p = 0.0015
\end{equation}

With $p = 0.0015 < 0.01$, the result is statistically significant, confirming that significant performance differences exist among the top 20 combinations.
This rejects the null hypothesis that ``all top combinations are equivalent and any can be used interchangeably,'' demonstrating that the choice of combination has a substantial impact on performance.

\textbf{Pathogen-Level BCa Bootstrap Confidence Intervals.}
\label{subsec:9pathogen_bootstrap}
The best single combination across all pathogens was evaluated using BCa bootstrap.
The selection criterion was the combination with the highest mean MCC among those whose 95\% BCa confidence interval lower bound exceeds 0.
Ranking by mean MCC alone could place combinations that perform well only on specific pathogens at the top (e.g., E\_only + SWAN achieves MCC = 0.311 on Norovirus but MCC = 0.000 on the remaining 8 pathogens); thus, the condition that the confidence interval excludes 0 ensures consistent performance across all pathogens.

\begin{itemize}
\item \textbf{Best single combination}: P*E + Wavelet ($t$ = 1.5)
\item \textbf{Mean MCC}: 0.124
\item \textbf{95\% BCa confidence interval}: [0.047, 0.202]
\end{itemize}

The mean MCC of the best single combination (0.124) is only 42\% of the mean oracle MCC (0.298), confirming that a single fixed combination falls substantially short. Despite appearing low, this MCC exceeds the random baseline (mean=0.001, SD=0.003) by approximately 41 standard deviations, and MCC $>$ 0.1 under extreme class imbalance (1--16\% positive rate) is not trivial.

\textbf{Subsampling Robustness.}
\label{subsec:9pathogen_subsampling}
The benchmark was repeated at MSA sequence proportions of 25\%, 50\%, 75\%, and 100\%, and the overall pathogen-average MCC was stable in the range of 0.054--0.058.
The MCC variation with decreasing sample size was within $\pm$0.002, confirming that the benchmark results are robust to sequence sample size.
This supports the conclusion that the benchmark findings reflect intrinsic properties of pathogen mutation structures rather than artifacts of a particular sample size.

%% ------------------------------------------------------------
\subsection{Supplementary Analyses}
\label{subsec:9pathogen_supplementary}

Several supplementary analyses were conducted to validate the robustness and generalizability of the benchmark findings.

\textbf{Cross-Paradigm Validation: Protein Language Model Scoring.}
\label{subsec:esm_validation}
To test whether the benchmark findings generalize beyond MSA-derived scoring formulas, an independent scoring paradigm was evaluated: the ESM-2 protein language model (650M parameters)~\cite{lin2023esm2}. ESM-2 computes per-position log-likelihood ratios (ESM-LLR) from a single consensus sequence using learned evolutionary representations, requiring no multiple sequence alignment. This provides orthogonal validation of the benchmark framework because the scoring signal originates from a fundamentally different methodology---a deep neural network trained on millions of protein families---rather than from the same MSA used by the nine existing scoring formulas.

The 650M-parameter ESM-2 model was selected as it represents the largest single-GPU model in the ESM-2 family, balancing computational feasibility with representation capacity. Larger models (ESM-2 3B) require multi-GPU inference, while smaller models (8M, 35M) may lack sufficient capacity for proteins exceeding 1,000 residues.
ESM-LLR was computed for all nine pathogens using wild-type marginal scoring: a single forward pass through the model yields $\text{ESM-LLR}[i] = -\log P(a_i \mid \text{context})$, where $a_i$ is the consensus amino acid at position $i$. For proteins exceeding the 1,022-token context window (SARS-CoV-2 Spike: 1,259 AA after MSA trimming; MERS Spike: 1,330 AA), a sliding window with 256-residue overlap was applied and scores were averaged in overlap regions. The resulting ESM-LLR arrays were evaluated through the same 39 detection methods used in the main benchmark.

\begin{table}[htbp]
\centering
\caption{Best MCC (Layer A, adaptive ground truth) for ESM-LLR versus the best existing scoring formula per pathogen. ESM rank indicates the position of ESM-LLR among 10 total scoring methods (9 existing + ESM-LLR) when each is paired with its best-performing detector.}
\label{tab:esm_comparison}
\small
\adjustbox{max width=\textwidth}{%
\begin{tabular}{lccccl}
\toprule
Pathogen & ESM-LLR & Existing Best & $\Delta$ & ESM Rank & Best Existing Score \\
\midrule
SARS-CoV-2 & 0.168 & 0.118 & +0.049 & 1/10 & E*rare \\
H3N2 & 0.360 & 0.549 & $-$0.189 & 5/10 & E\_only \\
Norovirus & 0.176 & 0.438 & $-$0.262 & 10/10 & P*E*rare \\
HIV-1 & 0.096 & 0.182 & $-$0.085 & 10/10 & E*rare \\
Dengue & 0.127 & 0.256 & $-$0.129 & 9/10 & E*rare \\
RSV & 0.196 & 0.249 & $-$0.053 & 5/10 & P*E \\
Influenza~B & 0.354 & 0.361 & $-$0.008 & 2/10 & H-score \\
MERS & 0.143 & 0.088 & +0.055 & 1/10 & rank(P*E) \\
HCV & 0.123 & 0.443 & $-$0.319 & 8/10 & E\_only \\
\midrule
Mean & 0.194 & 0.298 & $-$0.105 & 5.7/10 & --- \\
\bottomrule
\end{tabular}}% end adjustbox
\end{table}

Table~\ref{tab:esm_comparison} shows that ESM-LLR ranks first for two pathogens (SARS-CoV-2 and MERS) but performs poorly for others (last for Norovirus and HIV-1). The overall mean best MCC of ESM-LLR (0.194) is lower than the mean of existing scoring methods (0.298), indicating that MSA-derived scoring formulas generally provide more effective hotspot signals when sufficient sequence diversity is available. However, the pathogen-dependent rank variation (1st to 10th) reinforces the central finding of this benchmark: no single scoring paradigm dominates across all pathogens.

Spearman correlation between ESM-LLR and MSA-based scores varied substantially across pathogens. Norovirus showed strong positive correlation ($\rho = 0.63$ with E*rare, $p < 10^{-60}$), indicating that ESM-2's learned representations capture similar evolutionary signals as entropy-based scoring for this pathogen. In contrast, HIV-1 showed moderate negative correlation ($\rho = -0.35$, $p < 10^{-14}$), and Influenza~B showed near-zero correlation ($\rho = -0.04$, $p = 0.34$). For H3N2, ESM-LLR correlated positively with DMS fitness data ($\rho = 0.47$, $p < 10^{-30}$), suggesting that the protein language model captures experimentally measurable mutational tolerance. These results confirm that ESM-LLR provides a genuinely independent scoring signal for most pathogens, validating its use as a cross-paradigm check on the benchmark framework.
When ESM-LLR scores were additionally evaluated through all 39 detection methods (as a 10th scoring formula rather than cross-paradigm validation), the Friedman test across 10 scoring methods remained non-significant ($p = 0.608$), confirming that no single scoring paradigm---whether MSA-based or deep learning-based---consistently outperforms others across all pathogens.

\textbf{ANOVA Diagnostic Tests.}
\label{subsec:anova_diagnostics}
The Shapiro-Wilk test rejected residual normality ($W = 0.971$, $p = 3.29 \times 10^{-22}$), and Levene's test rejected homoscedasticity for all factors.
However, the non-normality is mild (skewness = 0.12, kurtosis = 2.42), and the large sample size ($n = 2{,}544$) ensures ANOVA robustness.
Kruskal-Wallis non-parametric tests confirmed all effects identically (pathogen: $H = 569.1$, $p < 10^{-117}$; detection family: $H = 86.1$, $p < 10^{-12}$; scoring: $H = 78.1$, $p < 10^{-13}$), maintaining the reliability of the ANOVA conclusions.

\textbf{Precision@$k$ Analysis.}
\label{subsec:precision_at_k}
The detection precision and recall of the best combination for each pathogen were analyzed. Table~\ref{tab:precision_at_k} presents the precision, recall, detection size, and best MCC for the top-1 combination of each pathogen.

\begin{table}[htbp]
\centering
\caption{Precision, recall, detection size, and best MCC of the top-1 combination for each pathogen.}
\label{tab:precision_at_k}
\adjustbox{max width=\textwidth}{
\begin{tabular}{lrrrr}
\toprule
Pathogen & Precision & Recall & $n_{\text{det}}$ & Best MCC \\
\midrule
Dengue       & 0.157 & 0.667 & 102 & 0.256 \\
H3N2         & 0.333 & 1.000 &  75 & 0.549 \\
HCV          & 0.622 & 0.426 &  37 & 0.443 \\
HIV-1        & 0.089 & 0.913 & 235 & 0.182 \\
Influenza B  & 0.189 & 0.824 &  74 & 0.361 \\
MERS         & 0.018 & 0.778 & 390 & 0.088 \\
Norovirus    & 0.290 & 0.733 &  38 & 0.438 \\
RSV          & 0.667 & 0.100 &   3 & 0.249 \\
SARS-CoV-2   & 0.090 & 0.240 &  67 & 0.118 \\
\midrule
Mean          & 0.273 & 0.631 &  -- & 0.298 \\
\bottomrule
\end{tabular}
}
\end{table}

The mean top-1 precision is 0.273 and recall is 0.631, with substantial heterogeneity across pathogens (e.g., H3N2 achieves perfect recall but RSV shows precision 0.667 with recall 0.10), confirming the limitations of single-threshold detection.

\textbf{DMS Validation.}
\label{subsec:dms_evolution_correlation}
\phantomsection\label{tab:dms_threshold_sensitivity}
For the 4 pathogens with DMS fitness data, Spearman correlations between DMS fitness and the best scoring formula output differed in direction (H3N2: $\rho = 0.191$; HIV-1: $\rho = 0.412$; RSV: $\rho = -0.363$; SARS-CoV-2: $\rho = -0.370$; all $p < 0.002$), confirming that DMS fitness and evolutionary adaptation capture different aspects of mutation impact and validating the 3-layer ground truth design.
Varying the DMS threshold from 10\% to 30\% altered the best-performing combination in all 4 pathogens (mean best DMS F1: 0.328 $\to$ 0.377 $\to$ 0.479), further supporting pathogen-dependent method selection.

%% ------------------------------------------------------------
\subsection{Summary of Key Findings}
\label{subsec:pahd_motivation_summary}

Synthesizing the key results of the 9-pathogen benchmark, the need for pathogen-adaptive method selection is demonstrated by four independent lines of evidence.

\begin{enumerate}
\item \textbf{Interaction dominance}: The detection family $\times$ pathogen interaction ($\omega^2 = 0.285$) exceeds all main effects. The best-performing method depends on the pathogen.
\item \textbf{Unique best combinations}: 9 unique best combinations were observed across 9/9 pathogens. No two pathogens share the same best combination.
\item \textbf{Generalization failure}: LOPO cross-validation yielded 0/9 combination matches with a mean generalization ratio of 4.9\%. The best combination selected from training pathogens does not generalize to new pathogens.
\item \textbf{Significant combination differences}: Friedman $\chi^2 = 42.44$, $p = 0.0015$. The choice of combination has a significant impact on performance.
\end{enumerate}

These results statistically establish that ``the best-performing method differs for each pathogen,'' motivating the conceptual PAHD framework---a future system that would adaptively select the best combination by profiling the mutation characteristics of each pathogen.

The following analysis applies the 9-pathogen benchmark results as a knowledge base for a proof-of-concept method selection experiment, demonstrating how MutBench's findings can be translated into practical method recommendations.

%% ------------------------------------------------------------
\subsection{PAHD Prototype Validation}
\label{subsec:pahd_prototype}

To test whether pathogens with similar mutation profiles share best-performing combinations, a proof-of-concept PAHD prototype was evaluated.
Two pathogen profile versions were compared: v1 (8 genomic features) and v2 (85 MSA-derived features including per-scoring statistics and spatial autocorrelation).
Using the best-performing Family-first strategy (predict detection family, then select within it), the v1 profile achieved 5/9 positive MCC (mean $= 0.005$) vs.\ the naive baseline (1/9 positive), and v2 improved mean MCC to 0.017 by reducing negative transfers (Table~\ref{tab:pahd_prototype}).

\begin{table}[htbp]
\centering
\caption{Per-pathogen test MCC: naive LOPO vs.\ PAHD Family-first prototype. Oracle is each pathogen's best combination MCC.}
\label{tab:pahd_prototype}
\adjustbox{max width=\textwidth}{
\begin{tabular}{lrrrr}
\toprule
\textbf{Pathogen} & \textbf{Oracle} & \textbf{Naive} & \textbf{PAHD Family v1} & \textbf{PAHD Family v2} \\
\midrule
Dengue      & 0.256 & 0.000 & $-$0.012 & $-$0.027 \\
H3N2        & 0.549 & 0.000 & 0.207 & 0.011 \\
HCV         & 0.443 & 0.000 & $-$0.291 & 0.141 \\
HIV-1       & 0.182 & 0.000 & 0.014 & 0.014 \\
Influenza B & 0.361 & 0.000 & $-$0.190 & $-$0.027 \\
MERS        & 0.088 & 0.000 & $-$0.012 & $-$0.025 \\
Norovirus   & 0.438 & 0.192 & 0.217 & 0.118 \\
RSV         & 0.249 & 0.000 & 0.069 & $-$0.054 \\
SARS-CoV-2  & 0.118 & 0.000 & 0.045 & 0.000 \\
\midrule
Mean         & 0.298 & 0.021 & 0.005 & 0.017 \\
Positive MCC count & 9/9 & 1/9 & \textbf{5/9} & 4/9 \\
\bottomrule
\end{tabular}
}
\end{table}

The exact combination match rate remains 0/9, confirming that an $n = 9$ knowledge base is insufficient for reliable combination-level prediction. However, at coarser granularities, family-level prediction recovered 48.3\% of oracle performance and scoring-level prediction 57.7\%, motivating the multi-feature integration approach developed in Chapter~\ref{ch:adapt}.

%% ============================================================
\section{Information-Type Analysis}
\label{sec:information_analysis}

The preceding sections demonstrated that the best-performing scoring--detection combination differs across pathogens ($\omega^2_{\text{family}\times\text{pathogen}} = 0.285$). This section investigates \textit{why} this interaction arises by analyzing the discriminative power of individual information types (features) that underlie the scoring formulas, and examines whether pathogen-specific feature importance patterns explain the observed cross-pathogen heterogeneity.

%% ------------------------------------------------------------
\subsection{Per-Feature Discriminative Power}
\label{subsec:per_feature_auc}

Ten positional features were evaluated as individual predictors of Layer~A (convergent evolution) ground truth across all 9 pathogens. These features represent six distinct information categories:

\begin{itemize}
\item \textbf{Frequency-based} (3 features): \textit{freq} (proportion of sequences carrying a non-reference amino acid at each position), \textit{entropy} (Shannon entropy measuring allelic diversity), and \textit{rare\_freq} (frequency of minor alleles below 5\% threshold, capturing low-frequency variants that may signal emerging adaptive mutations).
\item \textbf{Phylogeny-based} (1 feature): \textit{homoplasy} (number of independent mutation origins inferred by Fitch parsimony on the phylogenetic tree, distinguishing convergent evolution from founder-effect-driven frequency).
\item \textbf{Evolutionary rate} (1 feature): \textit{dN/dS proxy} (nonsynonymous substitution rate normalized by codon degeneracy of the consensus amino acid; since only protein MSAs were available, this approximates true dN/dS without nucleotide codon alignment).
\item \textbf{Deep learning-based} (2 features): \textit{esm2\_llr} (ESM-2 protein language model log-likelihood ratio~\cite{lin2023esm2,rives2021biological}, measuring how surprising a mutation is relative to learned evolutionary constraints from 250M protein sequences) and \textit{semantic change} (mean cosine distance between wild-type and mutant ESM-2 embeddings at each position~\cite{hie2021learning}, measuring how much a mutation alters the protein's functional context).
\item \textbf{Structure-based} (2 features): \textit{plddt} (AlphaFold2/ESMFold predicted local distance difference test score~\cite{jumper2021alphafold}, indicating structural confidence and disorder propensity) and \textit{SASA} (relative solvent accessible surface area computed from ESMFold-predicted structures using the Shrake--Rupley algorithm~\cite{shrake1973environment}, measuring surface exposure).
\item \textbf{Chemical property} (1 feature): \textit{Grantham distance} (mean physicochemical distance between the consensus and observed variant amino acids at each position~\cite{grantham1974amino}, capturing whether substitutions are conservative or radical).
\end{itemize}

These 10 features were selected from a broader set of 22 candidate information types identified from the literature. Features excluded from the analysis include: weighted contact number (WCN, high correlation with SASA, $r > 0.7$), intrinsic disorder prediction (largely redundant with pLDDT $< 50$~\cite{wilson2022alphafold}), co-mutation/mutual information (pairwise rather than per-position feature, incompatible with the univariate per-position analysis framework; additionally $O(L^2)$ computational cost), codon usage bias (weak association with hotspot status), and epitope/glycosylation distance (not uniformly available across all 9 pathogens). The selection criteria were: (1)~independence from existing features, (2)~applicability to all 9 pathogens, and (3)~prior validation in the literature.

Table~\ref{tab:per_feature_auc} and Figure~\ref{fig:feature_auc_heatmap} present the single-feature AUC (area under the ROC curve) for predicting Layer~A membership. Each feature was used as a standalone classifier via threshold sweeping (no model training), providing a direct measure of univariate discriminative power.

\begin{table}[htbp]
\centering
\caption{Per-feature AUC for Layer~A prediction across 9 pathogens (10 features). The highest AUC per pathogen is shown in \textbf{bold}. AUC = 0.5 indicates no discriminative power; AUC $< 0.5$ indicates inverse association.}
\label{tab:per_feature_auc}
\adjustbox{max width=\textwidth}{
\footnotesize
\begin{tabular}{lcccccccccc}
\toprule
\textbf{Pathogen} & \textbf{freq} & \textbf{entropy} & \textbf{rare} & \textbf{homoplasy} & \textbf{ESM-2} & \textbf{pLDDT} & \textbf{SASA} & \textbf{Grantham} & \textbf{dN/dS} & \textbf{Sem.~$\Delta$} \\
\midrule
SARS-CoV-2    & 0.522 & 0.521 & 0.522 & 0.491 & \textbf{0.586} & 0.513 & 0.504 & 0.522 & 0.523 & 0.523 \\
H3N2          & 0.548 & 0.533 & 0.535 & 0.639 & 0.405 & \textbf{0.787} & 0.416 & 0.554 & 0.548 & 0.649 \\
Norovirus     & 0.760 & 0.762 & 0.644 & \textbf{0.944} & 0.763 & 0.674 & 0.747 & 0.569 & 0.756 & 0.613 \\
HIV-1         & 0.713 & \textbf{0.730} & 0.592 & 0.686 & 0.339 & 0.359 & 0.533 & 0.561 & 0.627 & 0.645 \\
Dengue        & 0.577 & 0.577 & 0.548 & \textbf{0.770} & 0.521 & 0.678 & 0.474 & 0.603 & 0.558 & 0.495 \\
RSV           & 0.262 & 0.263 & 0.258 & 0.503 & 0.284 & \textbf{0.561} & 0.508 & 0.329 & 0.257 & 0.340 \\
Influenza~B   & 0.783 & 0.773 & 0.640 & \textbf{0.805} & 0.587 & 0.486 & 0.599 & 0.529 & 0.750 & 0.630 \\
MERS          & 0.515 & 0.544 & 0.664 & 0.681 & \textbf{0.695} & 0.162 & 0.579 & 0.470 & 0.477 & 0.665 \\
HCV           & 0.677 & 0.679 & 0.424 & 0.405 & 0.447 & 0.473 & 0.488 & 0.383 & \textbf{0.680} & 0.379 \\
\bottomrule
\end{tabular}
}
\begin{minipage}{\textwidth}
\smallskip
\footnotesize All AUC values computed on per-pathogen feature matrices using Layer~A (convergent evolution) as the binary label. $n$ ranges from 340 (HCV) to 1,330 (MERS) positions; Layer~A positive count ranges from 9 (MERS) to 54 (HCV). Column headers: rare = rare variant ratio, ESM-2 = ESM-2 LLR, dN/dS = dN/dS proxy, Sem.~$\Delta$ = semantic change.
\end{minipage}
\end{table}

\begin{figure}[htbp]
\centering
\includegraphics[width=\textwidth]{feature_auc_heatmap.png}
\caption{Per-feature AUC heatmap for Layer~A prediction. Green indicates high discriminative power; red indicates inverse association (AUC $< 0.5$). The best feature per pathogen is shown in bold. No single feature achieves the highest AUC for more than 3 of 9 pathogens.}
\label{fig:feature_auc_heatmap}
\end{figure}

The most discriminative feature varies across pathogens with no single feature dominating: homoplasy achieves the highest AUC for Norovirus (0.944), Dengue (0.770), and Influenza~B (0.805); pLDDT dominates for H3N2 (0.787) and RSV (0.561); ESM-2 LLR is best for SARS-CoV-2 (0.586) and MERS (0.695); entropy leads for HIV-1 (0.730); and dN/dS proxy marginally leads for HCV (0.680 vs.\ frequency 0.677). While the HCV difference is small (0.003), the expanded feature set provides measurable improvements in multivariate models (Section~\ref{subsec:rf_importance}).

Two pathogens exhibit inverse feature associations (AUC $< 0.5$): RSV shows AUC = 0.257--0.284 for frequency-based features, dN/dS proxy, and ESM-2 LLR, indicating that Layer~A positions in RSV tend to have \textit{lower} mutation frequency than non-hotspot positions. MERS exhibits a pLDDT AUC of 0.162, indicating that its convergent evolution positions are located in structurally disordered regions. These inverse associations further demonstrate that the biological relationship between features and hotspot status is pathogen-dependent.

%% ------------------------------------------------------------
\subsection{Feature Importance by Pathogen}
\label{subsec:rf_importance}

To assess feature importance in a multivariate context---where features may interact and provide complementary information---a Random Forest classifier~\cite{breiman2001random} (200 trees, max depth 5, balanced class weights, random state 42) was trained for each pathogen using the 10 features as predictors and Layer~A as the binary label. Feature importances were extracted as mean decrease in impurity (Gini importance), and model performance was evaluated via stratified 5-fold cross-validation AUC.

Figure~\ref{fig:feature_importance_grid} presents the feature importance distributions across all 9 pathogens. Three patterns are notable.

\textbf{First, the dominant feature differs across pathogens.}
For SARS-CoV-2, ESM-2 LLR dominates (importance = 0.331), followed by pLDDT (0.286) and the newly added SASA (0.202)---all three are non-frequency features. For H3N2, pLDDT remains the most important feature (0.338), with semantic change (0.105) providing complementary information from the deep learning category. For Norovirus, homoplasy dominates (0.345), consistent with its near-perfect univariate AUC (0.944). For HIV-1, ESM-2 LLR (0.214) and homoplasy (0.143) share importance. For HCV, the dN/dS proxy (0.145), homoplasy (0.139), and Grantham distance (0.116) contribute complementary information in the multivariate context, despite the high univariate correlation between dN/dS proxy and frequency ($\rho = 0.87$).

\textbf{Second, expanding from 6 to 10 features improves cross-validation AUC.}
Norovirus achieves the highest cross-validation AUC (0.899). SARS-CoV-2 improves from 0.554 (6 features) to 0.606 (10 features), and MERS improves from 0.537 to 0.875, demonstrating that the additional features---particularly SASA and semantic change---capture information relevant to hotspot prediction that was previously missing. RSV improves to 0.737, HCV to 0.701, and the remaining pathogens achieve: H3N2 (0.791), HIV-1 (0.802), Dengue (0.782), Influenza~B (0.756).

\textbf{Third, the importance rankings differ from univariate AUC rankings.}
For HIV-1, entropy has the highest univariate AUC (0.730) but ranks third in RF importance (0.127) behind ESM-2 LLR (0.214) and homoplasy (0.143), because entropy and frequency are highly correlated ($\rho = 0.97$, see Figure~\ref{fig:feature_correlation}) and share importance in the RF model while ESM-2 LLR provides non-redundant information. This demonstrates that multivariate analysis reveals complementary feature contributions that univariate analysis cannot capture.

\begin{figure}[htbp]
\centering
\includegraphics[width=\textwidth]{feature_importance_grid.png}
\caption{Random Forest feature importance for Layer~A prediction across 9 pathogens (10 features). Bar colors indicate feature category: blue (frequency-based), green (phylogeny/evolutionary), purple (deep learning), orange (structure), crimson (chemical property). Cross-validation AUC is shown in each panel title. The dominant feature differs across pathogens: ESM-2 LLR for SARS-CoV-2, pLDDT for H3N2, homoplasy for Norovirus, dN/dS proxy for HCV.}
\label{fig:feature_importance_grid}
\end{figure}

%% ------------------------------------------------------------
\subsection{Feature Correlation Structure}
\label{subsec:feature_correlation}

Understanding the correlation structure among the 10 features is essential for interpreting both the univariate and multivariate results. Figure~\ref{fig:feature_correlation} presents the mean pairwise Spearman correlation across all 9 pathogens.

\begin{figure}[htbp]
\centering
\includegraphics[width=0.85\textwidth]{feature_correlation_heatmap.png}
\caption{Mean pairwise Spearman correlation among 10 features across 9 pathogens. Three clusters are evident: frequency-derived features (freq, entropy, rare variant, dN/dS proxy) are highly correlated ($\rho > 0.54$); deep learning features (ESM-2 LLR, semantic change) show weak negative correlation ($\rho = -0.09$); and structure-based features (pLDDT, SASA) are weakly anti-correlated ($\rho = -0.19$). Homoplasy is the most independent feature ($|\rho| < 0.23$ with all others).}
\label{fig:feature_correlation}
\end{figure}

Three key patterns emerge. First, frequency, entropy, and dN/dS proxy form a highly correlated cluster ($\rho_{\text{freq--entropy}} = 0.97$, $\rho_{\text{freq--dN/dS}} = 0.87$), confirming that the protein-level dN/dS proxy shares substantial information with frequency-based features due to the absence of nucleotide-level synonymous substitution data (a limitation discussed in Section~\ref{sec:limitations}). Second, Grantham distance and semantic change show moderate correlation ($\rho = 0.52$), reflecting the relationship between physicochemical dissimilarity and embedding-space distance. Third, homoplasy is the most independent feature ($|\rho| < 0.23$ with all other features), explaining its unique discriminative power for pathogens where convergent evolution is a strong signal.

The correlation structure explains why expanding from 6 to 10 features yields differential improvements: SASA ($\rho = -0.19$ with pLDDT) provides genuinely independent structural information, while dN/dS proxy ($\rho = 0.87$ with frequency) adds limited independent signal. The largest improvements from the expanded set occur for pathogens where the new independent features (SASA, semantic change) are informative---notably MERS (CV AUC: 0.537 $\rightarrow$ 0.875) and SARS-CoV-2 (0.554 $\rightarrow$ 0.606).

%% ------------------------------------------------------------
\subsection{Implications for Adaptive Method Selection}
\label{subsec:info_type_implications}

The per-feature, Random Forest, and correlation analyses collectively provide a mechanistic explanation for the detection family $\times$ pathogen interaction ($\omega^2 = 0.285$) observed in the ANOVA (Section~\ref{subsec:9pathogen_statistical}). Different pathogens require different \textit{information types} for effective hotspot discrimination, and scoring formulas that emphasize the relevant information type for a given pathogen will produce higher-quality input signals for downstream detection algorithms.

This finding extends the PAHD framework from detection algorithm selection to information type selection. The 9-pathogen benchmark demonstrated that the best scoring--detection combination differs across pathogens (9/9 unique best combinations). The feature importance analysis reveals that this heterogeneity is rooted in pathogen-specific relationships between biological features and hotspot status: H3N2 hotspots are characterized by structural exposure (pLDDT), Norovirus hotspots by convergent evolution (homoplasy), SARS-CoV-2 hotspots by protein language model surprisal (ESM-2 LLR), and HCV hotspots by evolutionary rate (dN/dS proxy).

The oracle MCC ceiling (mean = 0.298 across 9 pathogens) observed in Section~\ref{subsec:9pathogen_bootstrap} reflects the discriminative capacity of the current scoring formulas, each of which relies on one or two of the 10 features. The Random Forest analysis demonstrates that combining all 10 features in a multivariate model achieves cross-validation AUC values up to 0.899 (Norovirus) and 0.875 (MERS), suggesting that the performance ceiling can be raised by integrating complementary information types into the scoring stage rather than relying on a single formula. This motivates a future direction for PAHD: pathogen-adaptive \textit{feature weighting} that selects and combines information types based on the pathogen's mutation profile, prior to applying a detection algorithm.

The implications of this feature heterogeneity for integrated detection are explored in Chapter~\ref{ch:adapt}, where multi-source integration analysis combines all 10 features into an equal-weight ensemble and evaluates it across 12 pathogens.
